\subsection{\texorpdfstring{ALGEBRAS OF TYPE \(D_{l}\) \((l\geq 2)\)}{ALGEBRAS OF TYPE D\_\{l\} (l\textbackslash geq 2)}}

\begin{enumerate}
\def\labelenumi{(\Roman{enumi})}
\tightlist
\item
  Let V be a vector space of even dimension \(2l \geq 4\) and let \(\varPsi\) be a non-degenerate symmetric bilinear form of maximum index l on V. By Algebra, Chap. IX, §4, no. 2, V can be written as the direct sum of two maximal totally isotropic subspaces F and \(F'\) . Let \((e_i)_{1 \leq i \leq l}\) be a basis of F and \((e_{-i})_{1 \leq i \leq l}\) the dual basis of \(F'\) (for the duality between F and \(F'\) defined by \(\varPsi\) ). Then \(e_1, \ldots, e_l, e_{-l}, \ldots, e_{-1}\) is a basis of V; we shall call it a Witt basis of V. The matrix of \(\varPsi\) with respect to this basis is the square matrix S of order 2l all of whose entries are zero, except those situated on the second diagonal, which are equal to 1. The algebra \(\mathfrak{g} = \mathfrak{o}(\varPsi)\) can be identified with the algebra \(\mathfrak{o}_S(2l, k)\) of square matrices g of order 2l such that \(g = -S^t g S\) . It has dimension \(l(2l - 1)\) . An easy calculation shows that g is the set of matrices of the form
\end{enumerate}

\[
\left( \begin{array}{c c} A & B \\ C & D \end{array} \right)
\]

where A, B, C, D are square matrices of order l such that \(B = -s^{t}Bs\) , \(C = -s^{t}Cs\) and \(D = -s^{t}As\) (s is the matrix of order l all of whose entries are zero except those situated on the second diagonal which are equal to 1).

Let \(\mathfrak{h}\) be the set of diagonal matrices belonging to \(\mathfrak{g}\) . This is a commutative subalgebra of \(\mathfrak{g}\) , with basis formed by the elements \(H_{i} = E_{i,i} - E_{-i,-i}\) for \(1 \leq i \leq l\) . Let \((\varepsilon_{i})\) be the basis of \(\mathfrak{h}^{*}\) dual to \((H_{i})\) . Put, for \(1 \leq i < j \leq l\) ,

\[
\left\{ \begin{array}{l l} X _ {\varepsilon_ {i} - \varepsilon_ {j}} & = E _ {i, j} - E _ {- j, - i} \\ X _ {- \varepsilon_ {i} + \varepsilon_ {j}} & = - E _ {j, i} + E _ {- i, - j} \\ X _ {\varepsilon_ {i} + \varepsilon_ {j}} & = E _ {i, - j} - E _ {j, - i} \\ X _ {- \varepsilon_ {i} - \varepsilon_ {j}} & = - E _ {- j, i} + E _ {- i, j}. \end{array} \right.\tag{11}
\]

These elements form a basis of a complement of h in g. For \(h \in h\) ,

\[
[ h, X _ {\alpha} ] = \alpha (h) X _ {\alpha}
\]

for all \(\alpha \in \mathbb{R}\) , where \(\mathbb{R}\) is the set of the \(\pm \varepsilon_i \pm \varepsilon_j\) ( \(i < j\) ). Thus, \(\mathfrak{h}\) is a splitting Cartan subalgebra of \(\mathfrak{g}\) , and the roots of \((\mathfrak{g}, \mathfrak{h})\) are the elements of \(\mathbb{R}\) . The root system \(\mathbb{R}\) of \((\mathfrak{g}, \mathfrak{h})\) is thus of type \(\mathrm{D}_l\) for \(l \geq 3\) , of type \(\mathrm{A}_1 \times \mathrm{A}_1\) (in other words of type \(\mathrm{D}_2\) ) for \(l = 2\) (Chap. VI, §4, no. 8.I extended to the case \(l = 2\) ). Consequently, \(\mathfrak{g}\) is a splittable simple Lie algebra of type \(\mathrm{D}_l\) if \(l \geq 3\) .

Every splitting Cartan subalgebra of g is transformed into h by an elementary automorphism of g, and hence by an element of \(\mathbf{O}(\varPsi)\) (cf.~(VII)) and consequently is the set \(h_{\beta}\) of elements of g whose matrix with respect to a Witt basis \(\beta\) of V is diagonal. We verify immediately that the only subspaces invariant under \(h_{\beta}\) are those generated by a subset of \(\beta\) .

Since the algebras \(\mathfrak{o}_{S}(4,k)\) and \(\mathfrak{sl}(2,k)\times\mathfrak{sl}(2,k)\) have the same root systems, they are isomorphic. Similarly, \(\mathfrak{o}_{S}(6,k)\) and \(\mathfrak{sl}(4,k)\) are isomorphic (cf.~also Exerc. 3). From now on, we assume that \(l\geq3\) .

\begin{enumerate}
\def\labelenumi{(\Roman{enumi})}
\setcounter{enumi}{1}
\tightlist
\item
  We determine \(\mathbb{R}^{\vee}\) by means of Chap. VI, §4, no. 8.V. We find that
\end{enumerate}

\[
H _ {\varepsilon_ {i} - \varepsilon_ {j}} = H _ {i} - H _ {j}, \quad H _ {\varepsilon_ {i} + \varepsilon_ {j}} = H _ {i} + H _ {j}.
\]

\begin{enumerate}
\def\labelenumi{(\Roman{enumi})}
\setcounter{enumi}{2}
\tightlist
\item
  Put \(\alpha_{1} = \varepsilon_{1} - \varepsilon_{2},\alpha_{2} = \varepsilon_{2} - \varepsilon_{3},\ldots ,\alpha_{l - 1} = \varepsilon_{l - 1} - \varepsilon_{l},\alpha_{l} = \varepsilon_{l - 1} + \varepsilon_{l}\) . By Chap. VI, §4, no. 8.II, \((\alpha_{1},\dots,\alpha_{l})\) is a basis B of R; the positive roots relative to B are the \(\varepsilon_i\pm \varepsilon_j\) ( \(i < j\) ). The corresponding Borel subalgebra \(\mathfrak{b}\) is the set of upper triangular matrices belonging to \(\mathfrak{g}\) .
\end{enumerate}

It is easily verified that the only non-trivial vector subspaces invariant under b are the totally isotropic subspaces \(V_{1},\ldots,V_{l},V_{l}^{\prime}\) , where \(V_{i}\) is generated by \(e_{1},\ldots,e_{i}\) and \(V_{l}^{\prime}\) by \(e_{1},\ldots,e_{l-1},e_{-l}\) , and the orthogonal complements \(V_{-1},\ldots,V_{-l+1}\) of \(V_{1},\ldots,V_{l-1}\) ; the orthogonal complement \(V_{-i}\) of \(V_{i}\) is generated by \(e_{1},\ldots,e_{l},e_{-l},\ldots,e_{-(i+1)}\) . But an immediate calculation shows that, if an element \(a\in g\) leaves \(V_{l-1}\) stable, its matrix is of the form

\[
\left( \begin{array}{c c c} A & x & B \\ 0 & \left( \begin{array}{c c} \lambda & 0 \\ 0 & - \lambda \end{array} \right) & y \\ 0 & 0 & D \end{array} \right)
\]

where \(A, B, D\) are square matrices of order \(l - 1\) , \(x\) (resp. \(y\) ) is a matrix with 2 columns and \(l - 1\) rows (resp. 2 rows and \(l - 1\) columns), and \(\lambda \in k\) . It follows that \(a\) leaves \(\mathrm{V}_l\) and \(\mathrm{V}_l'\) stable. Consequently, \(\mathfrak{b}\) is the set of \(a \in \mathfrak{g}\) leaving all the elements of the isotropic flag \((\mathrm{V}_1, \ldots, \mathrm{V}_{l-1})\) stable. Note that the preceding and Witt's theorem (Algebra, Chap. IX, §4, no. 3, Th. 1) imply that \(\mathrm{V}_l\) and \(\mathrm{V}_l'\) are the only maximal totally isotropic subspaces containing \(\mathrm{V}_{l-1}\) .

We say that an isotropic flag is quasi-maximal if it is composed of l-1 totally isotropic subspaces of dimensions \(1,\ldots,l-1\) . We then see as in no. 2 that, for any quasi-maximal isotropic flag \(\delta\) , the set \(b_{\delta}\) of \(a\in g\) leaving the elements of \(\delta\) stable is a Borel subalgebra of g and that the map \(\delta\mapsto b_{\delta}\) is a bijection from the set of quasi-maximal isotropic flags to the set of Borel subalgebras.

We say that an isotropic flag is proper if it does not contain both a subspace of dimension \(l\) and a subspace of dimension \(l-1\) . Let \(\delta\) be such an isotropic flag and let \(\mathfrak{p}_{\delta}\) be the set of \(a \in \mathfrak{g}\) leaving stable the elements of \(\delta\) . If \(\delta \subset \{\mathrm{V}_1, \ldots, \mathrm{V}_l, \mathrm{V}_l'\}\) , then \(\mathfrak{p}_{\delta}\) is a parabolic subalgebra of \(\mathfrak{g}\) , containing \(\mathfrak{b}\) , and it is easy to verify that the only totally isotropic subspaces \(\neq \{0\}\) stable under \(\mathfrak{p}_{\delta}\) are the elements of \(\delta\) . Since there are \(2^{l-2}\) proper isotropic flags contained in \(\{\mathrm{V}_1, \ldots, \mathrm{V}_l, \mathrm{V}_l'\}\) and containing \(\mathrm{V}_{l-1}\) (resp. \(\mathrm{V}_l\) , resp. \(\mathrm{V}_l'\) , resp. containing neither \(\mathrm{V}_{l-1}\) , nor \(\mathrm{V}_l\) , nor \(\mathrm{V}_l'\) ), this gives \(2^l\) parabolic subalgebras containing \(\mathfrak{b}\) . It follows as above that the map \(\delta \mapsto \mathfrak{p}_{\delta}\) is a bijection from the set of proper isotropic flags to the set of parabolic subalgebras of \(\mathfrak{g}\) .

\begin{enumerate}
\def\labelenumi{(\Roman{enumi})}
\setcounter{enumi}{3}
\tightlist
\item
  The fundamental weights corresponding to \(\alpha_{1},\ldots,\alpha_{l}\) are, by Chap. VI, §4, no. 8.VI,
\end{enumerate}

\[
\varpi_ {i} = \varepsilon_ {1} + \varepsilon_ {2} + \dots + \varepsilon_ {i} (1 \leq i \leq l - 2)
\]

\[
\begin{array}{c} \varpi_ {l - 1} = \frac {1}{2} (\varepsilon_ {1} + \varepsilon_ {2} + \ldots + \varepsilon_ {l - 2} + \varepsilon_ {l - 1} - \varepsilon_ {l}) \\ \varpi_ {l} = \frac {1}{2} (\varepsilon_ {1} + \varepsilon_ {2} + \ldots + \varepsilon_ {l - 2} + \varepsilon_ {l - 1} + \varepsilon_ {l}). \end{array}
\]

Let \(\sigma\) be the identity representation of \(\mathfrak{g}\) on V. The exterior power \(\bigwedge^{r}\sigma\) operates on \(\mathrm{E} = \bigwedge^{r}(\mathrm{V})\) . If \(h\in \mathfrak{h}\) , we have

\[
\sigma (h) e _ {i} = \varepsilon_ {i} (h) e _ {i}, \quad \sigma (h) e _ {- i} = - \varepsilon_ {i} (h) e _ {- i}
\]

for \(1 \leq i \leq l\) . It follows that, for \(1 \leq r \leq l\) , \(\varepsilon_{1} + \cdots + \varepsilon_{r}\) is the highest weight of \(\bigwedge^{r} \sigma\) , the elements of weight \(\varepsilon_{1} + \cdots + \varepsilon_{r}\) being those proportional to \(e_{1} \wedge \cdots \wedge e_{r}\) .

We shall show that, for \(1 \leq r \leq l - 2\) , the representation \(\bigwedge^{r} \sigma\) is a fundamental representation of weight \(\varpi_{r}\) .

For this, it suffices to show that \(\bigwedge^{r}\sigma\) is irreducible for \(1 \leq r \leq l - 1\) (note that the representation \(\bigwedge^{l}\sigma\) is not irreducible, cf.~Exerc. 10), or that the smallest subspace \(T_{r}\) of \(\bigwedge^{r}V\) containing \(e_{1} \wedge \cdots \wedge e_{r}\) and stable under g is the whole of \(\bigwedge^{r}V\) . This is immediate for r = 1. For r = 2, we see as in no. 2 that \(\bigwedge^{2}\sigma\) is equivalent to the adjoint representation of g, which is irreducible since g is simple. The proof is completed by arguing by induction on l, as in no. 2, but assuming that \(l - 1 \geq r \geq 3\) .

We are now going to determine the fundamental representations of highest weight \(\varpi_{l-1}\) and \(\varpi_{l}\) . Let Q be the quadratic form \(x \mapsto \frac{1}{2}\Psi(x, x)\) . We have defined in no. 2.IV the spinor representation \(\lambda\) of the Clifford algebra \(\mathrm{C}(\mathrm{Q})\) on \(N = \bigwedge F'\) . We verify immediately that the subspace \(N_{+}\) (resp. \(N_{-}\) ) of N given by the sum of the \(\bigwedge^{p}F'\) for p even (resp. odd) is stable under the restriction of \(\lambda\) to \(\mathrm{C}^{+}(\mathrm{Q})\) . Consequently, the representations \(\lambda_{+}\) and \(\lambda_{-}\) of \(\mathrm{C}^{+}(\mathrm{Q})\) on \(N_{+}\) and \(N_{-}\) respectively are the semi-spinor representations of \(\mathrm{C}^{+}(\mathrm{Q})\) (Algebra, Chap. IX, §9, no. 4); they are irreducible, of dimension \(2^{l-1}\) and inequivalent. Let \(\rho_{+} = \lambda_{+} \circ f\) and \(\rho_{-} = \lambda_{-} \circ f\) be the corresponding irreducible representations of g (no. 2, Lemma 1 (vi)). In view of Lemma 1 (i), we have

\[
f (H _ {i}) = \frac {1}{2} (e _ {i} e _ {- i} - e _ {- i} e _ {i}) = e _ {i} e _ {- i} - \frac {1}{2} = \frac {1}{2} - e _ {- i} e _ {i}
\]

and we see, as in no. 2.IV, that, for \(h \in h\) and \(1 \leq i_{1} < \cdots < i_{k} \leq l\) ,

\[
\begin{array}{l} \lambda \circ f (h) (e _ {- i _ {1}} \wedge \dots \wedge e _ {- i _ {k}}) \\ = (\frac {1}{2} (\varepsilon_ {1} + \dots + \varepsilon_ {l}) - (\varepsilon_ {i _ {1}} + \dots + \varepsilon_ {i _ {k}})) (h) (e _ {- i _ {1}} \wedge \dots \wedge e _ {- i _ {k}}). \end{array}
\]

Consequently, the highest weight of \(\rho_{+}\) (resp. \(\rho_{-}\) ) is \(\varpi_{l}\) (resp. \(\varpi_{l-1}\) ).

We say that \(\rho_{+}\) and \(\rho_{-}\) are the semi-spinor representations of g. All their weights are simple. We also say that \(\rho = \lambda \circ f = \rho_{+} \oplus \rho_{-}\) is the spinor representation of g.

\begin{enumerate}
\def\labelenumi{(\Alph{enumi})}
\setcounter{enumi}{21}
\tightlist
\item
  For \(1 \leq r \leq l - 2\) , the fundamental representation \(\bigwedge^{r} \sigma\) is orthogonal: it leaves invariant the extension of \(\varPsi\) to \(\bigwedge^{r} \mathrm{V}\) .
\end{enumerate}

Consider now the spinor representation \(\rho\) of \(\mathfrak{g}\) . We show as in no. 2 that, for \(1 \leq i, j \leq l, i \neq j\) ,

\[
\begin{array}{r} f (X _ {\varepsilon_ {i} - \varepsilon_ {j}}) = \pm e _ {i} e _ {- j} \\ f (X _ {\varepsilon_ {i} + \varepsilon_ {j}}) = \pm e _ {i} e _ {j} \\ f (X _ {- \varepsilon_ {i} - \varepsilon_ {j}}) = \pm e _ {- i} e _ {- j}. \end{array}
\]

It follows that the non-degenerate bilinear form \(\varPhi\) introduced in no. 2.V is invariant under \(\rho(\mathfrak{g})\) . Thus, the spinor representation \(\rho\) leaves invariant a non-degenerate form that is symmetric for \(l \equiv 0, -1 \pmod{4}\) and alternating for \(l \equiv 1, 2 \pmod{4}\) .

If \(l\) is even, the restrictions of \(\varPhi\) to \(\mathrm{N}_{+}\) and \(\mathrm{N}_{-}\) are non-degenerate and the semi-spinor representations are orthogonal for \(l \equiv 0 \pmod{4}\) and symplectic for \(l \equiv 2 \pmod{4}\) . Moreover, we remark that \(w_{0} = -1\) (Chap. VI, §4, no. 8.XI).

On the other hand, if \(l\) is odd, \(\mathrm{N}_{+}\) and \(\mathrm{N}_{-}\) are totally isotropic for \(\varPhi\) . Moreover, \(-w_{0}(\alpha_{i}) = \alpha_{i}\) for \(1 \leq i \leq l - 2\) , \(-w_{0}(\alpha_{l}) = \alpha_{l - 1}\) and \(-w_{0}(\alpha_{l - 1}) = \alpha_{l}\) (Chap. VI, §4, no. 8.XI), so \(-w_{0}(\varpi_{l}) = \varpi_{l - 1}\) and the semi-spinor representations are neither orthogonal nor symplectic; each of them is isomorphic to the dual of the other.

\begin{enumerate}
\def\labelenumi{(\Roman{enumi})}
\setcounter{enumi}{5}
\tightlist
\item
  For all \(x \in \mathfrak{g}\) , the characteristic polynomial of \(\sigma(x)\) takes the form
\end{enumerate}

\[
\mathrm{T} ^ {2 l} + f _ {1} (x) \mathrm{T} ^ {2 l - 1} + \dots + f _ {2 l} (x).
\]

We see as in no. 3 that \(f_{1}=f_{3}=f_{5}=\cdots=0\) . By Chap. VI, §4, no. 8.IX and §8, no. 3, Th. 1, there exists a polynomial function \(\tilde{f}\) on g such that \(f_{2}, f_{4}, \ldots, f_{2l-2}, \tilde{f}\) generate the algebra \(\mathrm{I}(\mathfrak{g}^{*})\) of invariant polynomial functions on g, are algebraically independent, and further \(\tilde{f}^{2}=(-1)^{l}f_{2l}\) .

For all \(x \in \mathfrak{g}\) , we have \({}^t(Sx) = {}^t xS = -Sx\) , so we can consider \(\operatorname{Pf}(Sx)\) , which is a polynomial function of \(x\) . Now:

\[
f _ {2 l} (x) = \det (x) = (- 1) ^ {l} \det (S x) = (- 1) ^ {l} (\operatorname{Pf} (S x)) ^ {2}.
\]

Thus, we can take \(\tilde{f}(x) = \mathrm{Pf}(Sx)\) .

\begin{enumerate}
\def\labelenumi{(\Roman{enumi})}
\setcounter{enumi}{6}
\tightlist
\item
  Recall (§5, no. 3, Cor. 1 of Prop. 5) that \(\mathrm{Aut}(\mathfrak{g}) / \mathrm{Aut}_0(\mathfrak{g})\) can be identified with the group \(\mathrm{Aut}(\mathrm{D})\) of automorphisms of the Dynkin graph D of \((\mathfrak{g},\mathfrak{h})\) . When \(l\neq 4\) , \(\mathrm{Aut}(\mathrm{D})\) is the group of order 2 consisting of the permutations of \(\alpha_{1},\ldots ,\alpha_{l}\) that leave \(\alpha_{1},\ldots ,\alpha_{l - 2}\) fixed. When \(l = 4\) , \(\mathrm{Aut}(\mathrm{D})\) consists of the permutations of \(\alpha_{1},\ldots ,\alpha_{4}\) that leave \(\alpha_{2}\) fixed; it is isomorphic to \(\mathfrak{S}_3\) (cf.~Chap. VI, §4, no. 8.XI). In all cases, the subgroup of \(\mathrm{Aut}(\mathrm{D})\) consisting of the elements that leave \(\alpha_{1}\) fixed is of order 2. We denote by \(\operatorname{Aut}'(\mathfrak{g})\) the corresponding subgroup of \(\operatorname{Aut}(\mathfrak{g})\) ; we have \(\operatorname{Aut}'(\mathfrak{g}) = \operatorname{Aut}(\mathfrak{g})\) if \(l \neq 4\) and \((\operatorname{Aut}(\mathfrak{g}) : \operatorname{Aut}'(\mathfrak{g})) = 3\) if l = 4; moreover,
\end{enumerate}

\[
\left(\operatorname{Aut} ^ {\prime} (\mathfrak {g}): \operatorname{Aut} _ {0} (\mathfrak {g})\right) = 2.
\]

An element \(s \in \operatorname{Aut}(\mathfrak{g})\) belongs to \(\operatorname{Aut}'(\mathfrak{g})\) if and only if \(\sigma \circ s\) is equivalent to \(\sigma\) (this follows from the fact that \(\varpi_{1}\) is the highest weight of \(\sigma\) ). We conclude as in no. 2.VII that \(\operatorname{Aut}'(\mathfrak{g})\) can be identified with \(\Sigma/k^{*}\) , where \(\Sigma\) is the group of similarities of V relative to \(\Psi\) .

Let \(s \in \Sigma\) , and let \(\lambda(s)\) be the multiplier of \(s\) . We have \(\det(s) = \lambda(s)^l\) if \(s\) is direct, and \(\det(s) = -\lambda(s)^l\) if \(s\) is inverse (Algebra, Chap. IX, §6, no. 5). The direct similarities form a subgroup \(\Sigma_0\) of index 2 in \(\Sigma\) ; we have \(\Sigma_0 \supset k^*\) . The group \(\Sigma_0 / k^*\) is equal to the subgroup \(\operatorname{Aut}_0(\mathfrak{g})\) of \(\operatorname{Aut}'(\mathfrak{g}) = \Sigma / k^*\) . Indeed, it suffices to verify this when \(k\) is algebraically closed: in that case \(\operatorname{Aut}_0(\mathfrak{g}) = \operatorname{Aut}_e(\mathfrak{g})\) is equal to its derived group (§11, no. 2, Prop. 3), hence is contained in \(\Sigma_0 / k^*\) , and since they are both of index 2 in \(\Sigma / k^*\) , they are equal.

On the other hand, as in no. 3.VII there is an exact sequence

\[
1 \longrightarrow \mathbf {S O} (\varPsi) / \{1, - 1 \} \longrightarrow \varSigma_ {0} / k ^ {*} \longrightarrow k ^ {*} / k ^ {* 2} \longrightarrow 1.
\]

Identify \(\mathbf{SO}(\varPsi)/\{1,-1\}\) with a subgroup of \(\Sigma_{0}/k^{*}=\mathrm{Aut}_{0}(\mathfrak{g})\) . Since \(k^{*}/k^{*2}\) is commutative, we have \(\mathrm{Aut}_{e}(\mathfrak{g})\subset\mathbf{SO}(\varPsi)/\{1,-1\}\) . In fact, it can be shown (Exerc. 11) that \(\mathrm{Aut}_{e}(\mathfrak{g})\) is equal to the image in \(\mathbf{SO}(\varPsi)/\{1,-1\}\) of the reduced orthogonal group \(\mathbf{O}_{0}^{+}(\varPsi)\) of \(\varPsi\) (Algebra, Chap. IX, §9, no. 5).

\begin{enumerate}
\def\labelenumi{(\Roman{enumi})}
\setcounter{enumi}{7}
\tightlist
\item
  The canonical bilinear form \(\varPhi_{\mathrm{R}}\) on \(\mathfrak{h}^*\) is given by
\end{enumerate}

\[
\Phi_ {\mathrm{R}} \left(\xi_ {1} \varepsilon_ {1} + \dots + \xi_ {l} \varepsilon_ {l}, \xi_ {1} ^ {\prime} \varepsilon_ {1} + \dots + \xi_ {l} ^ {\prime} \varepsilon_ {l}\right) = \frac {1}{4 (l - 1)} \left(\xi_ {1} \xi_ {1} ^ {\prime} + \dots + \xi_ {l} \xi_ {l} ^ {\prime}\right)
\]

(Chap. VI, §4, no. 8.V). Thus, the restriction of the Killing form to \(\mathfrak{h}\) is

\[
\Phi \left(\xi_ {1} H _ {1} + \dots + \xi_ {l} H _ {l}, \xi_ {1} ^ {\prime} H _ {1} + \dots + \xi_ {l} ^ {\prime} H _ {l}\right) = 4 (l - 1) \left(\xi_ {1} \xi_ {1} ^ {\prime} + \dots + \xi_ {l} \xi_ {l} ^ {\prime}\right).
\]

\begin{enumerate}
\def\labelenumi{(\Roman{enumi})}
\setcounter{enumi}{8}
\tightlist
\item
  Recall the \(X_{\alpha}\) ( \(\alpha \in \mathbb{R}\) ) defined by formulas (11). We verify easily that \([X_{\alpha}, X_{-\alpha}] = -H_{\alpha}\) for \(\alpha \in \mathbb{R}\) . On the other hand, the map \(\theta : a \mapsto -^{t}a\) is an automorphism of \(\mathfrak{g}\) and \(\theta(X_{\alpha}) = X_{-\alpha}\) for all \(\alpha \in \mathbb{R}\) . Consequently \((X_{\alpha})_{\alpha \in \mathbb{R}}\) is a Chevalley system in \((\mathfrak{g}, \mathfrak{h})\) .
\end{enumerate}

Assume that \(k = \mathbf{Q}\) . By Chap. VI, §4, no. 8.VIII, the subalgebra \(\mathfrak{h}\) has three permissible lattices if \(l\) is odd and four permissible lattices if \(l\) is even. In particular, the lattice \(\mathcal{H}\) generated by the \(H_i\) is permissible. But this lattice is the set of diagonal matrices in \(\mathfrak{g}\) with integer entries. It follows that \(\mathfrak{o}_S(2l,\mathbf{Z})\) is the Chevalley order \(\mathcal{H} + \sum \mathbf{Z}.X_\alpha\) in \(\mathfrak{g}\) . Since \(X_{\alpha}^{2} = 0\) for all \(\alpha \in \mathrm{R}\) , we see that the lattice \(\mathcal{V}\) in \(\mathrm{V}\) generated by the Witt basis \((e_i)\) is an admissible lattice in \(\mathrm{V}\) for \(\mathfrak{o}_S(2l,\mathbf{Z})\) . The same holds for \(\bigwedge^r\mathcal{V}\) in \(\bigwedge^r\mathrm{V}\) .

On the other hand, if we take \(\mathrm{P}(\mathrm{R}^{\vee}) = \mathbf{Z}\) . \(\frac{1}{2} \sum_{i=1}^{l} H_i + \mathcal{H}\) as permissible lattice and \(\mathcal{G} = \mathrm{P}(\mathrm{R}^{\vee}) + \sum \mathbf{Z}. X_{\alpha}\) as Chevalley order, we see that \(\bigwedge^r \mathrm{V}\) has an admissible lattice only if \(r\) is even; \(\bigwedge^r \mathcal{V}\) is then admissible.

Consider the reductive Lie algebra \(\mathfrak{o}(\varPsi) + \mathbf{Q}.1\) ; we see immediately that the lattice \(\tilde{\mathcal{G}} = (\mathfrak{o}(\varPsi) + \mathbf{Q}.1) \cap \mathfrak{gl}(2l, \mathbf{Z})\) is a Chevalley order. The Chevalley order \(\mathcal{G}\) is the projection of \(\tilde{\mathcal{G}}\) onto \(\mathfrak{o}(\varPsi)\) parallel to the centre \(\mathbf{Q}.1\) .

Finally, we see as in no. 2 that the lattice \(N_{+}\) (resp. \(N_{-}\) ) generated by the \(e_{-i_{1}}\wedge\cdots\wedge e_{-i_{2k}}\) (resp. \(e_{-i_{1}}\wedge\cdots\wedge e_{-i_{2k+1}}\) ) is admissible for the semi-spinor representation of the Chevalley order \(\mathrm{Q}(\mathrm{R}^{\vee})+\sum_{\alpha\in\mathbb{R}}\mathbf{Z}.X_{\alpha}\) . On the other hand, \(N_{+}\) and \(N_{-}\) have no admissible lattice for \(\mathfrak{o}_{S}(2l,\mathbf{Z})\) .

